\chapter{Introducción}
\label{chap:introduccion}

El electrocardiograma es una señal temporal, pero una parte importante de su circulación clínica ocurre como imagen: papel, PDF, captura de pantalla o trazado incrustado en una historia. Esta doble naturaleza abre una pregunta práctica: si el dato disponible ya es visual, ¿merece la pena entrenar un modelo de visión específico o puede aprovecharse un modelo multimodal general mediante ejemplos en contexto?

Este Trabajo de Fin de Grado aborda esa pregunta en un escenario deliberadamente difícil: pocos datos etiquetados, trazados electrocardiográficos y sospecha de síndrome de Brugada. El objetivo no es construir un producto sanitario ni un clasificador diagnóstico autónomo. La orientación correcta es la de triaje clínico: priorizar trazados que merecen revisión temprana por especialistas y generar evidencia reproducible sobre qué familia de modelos resulta más útil en cada régimen de datos.

\section{Motivación clínica}

El síndrome de Brugada se asocia a un patrón electrocardiográfico característico en derivaciones precordiales derechas, especialmente el patrón tipo 1. Sin embargo, el diagnóstico no se reduce a reconocer una forma de onda. El ECG puede ser dinámico, el patrón puede aparecer de forma espontánea, durante fiebre o tras una prueba con bloqueantes del canal de sodio, y existen fenocopias que deben excluirse \cite{zeppenfeld2022,wilde2023scbt}.

Por esta razón, el sistema propuesto no responde a la pregunta "¿tiene este paciente Brugada?". Responde a una pregunta previa y más prudente: "¿debe este trazado revisarse antes y, si procede, continuar el circuito clínico?". En un flujo asistencial real, una alerta automática tendría que terminar en revisión experta, posible registro con derivaciones V1-V2 en espacios intercostales superiores y, solo en pacientes seleccionados, una prueba farmacológica monitorizada. La decisión final seguiría siendo clínica.

\begin{figure}[H]
  \centering
  \safeincludegraphics[width=0.96\textwidth]{results/fig_triaje_clinico.pdf}
  \caption{Uso previsto del sistema: triaje de sospecha y priorización de revisión, no diagnóstico autónomo.}
  \label{fig:triaje_intro}
\end{figure}

\section{Motivación metodológica}

Los modelos supervisados de ECG funcionan bien cuando existen cohortes grandes, diversas y bien etiquetadas. Trabajos previos han mostrado que las redes convolucionales pueden extraer información diagnóstica de trazados renderizados a gran escala \cite{sangha2022}. El problema aparece al comienzo de un proyecto biomédico, cuando solo hay pocos ejemplos revisados y la enfermedad o la definición operativa son raras.

Este TFG compara dos formas de usar esos pocos ejemplos. La primera es el enfoque clásico: una CNN usa los \(K\) ejemplos para actualizar sus pesos. La segunda es el enfoque de aprendizaje en contexto: un modelo de visión y lenguaje mantiene sus pesos fijos y recibe los \(K\) ejemplos dentro del mensaje. La pregunta de investigación no es cuál tecnología es más moderna, sino a partir de qué punto compensa entrenar y cuándo resulta más razonable aprovechar conocimiento visual previo mediante ICL.

La fase cerrada de esta memoria corresponde a la primera rama. Se entrena una CNN ResNet18, se evalúa en un simulador controlado, se traslada a Brugada HUCA y se revisa una rama de adaptación de dominio. Esa evidencia no pretende cerrar la comparación contra VLM. Pretende establecer la línea base que la siguiente fase debe superar o matizar.

\begin{figure}[H]
  \centering
  \safeincludegraphics[width=0.96\textwidth]{results/fig_flujo_conceptual.pdf}
  \caption{Flujo conceptual del TFG: misma entrada visual, dos formas de usar pocos ejemplos y lectura clínica común de triaje.}
  \label{fig:flujo_conceptual}
\end{figure}

\section{Pregunta de investigación}

La pregunta general del trabajo es:

\begin{quote}
En ECG como imagen y con pocos datos etiquetados, ¿cuándo merece la pena entrenar una CNN específica y cuándo puede ser preferible usar modelos de visión y lenguaje con aprendizaje en contexto para apoyar triaje de sospecha de Brugada?
\end{quote}

La memoria actual responde de forma completa a la parte CNN de esa pregunta mediante tres subpreguntas:

\begin{enumerate}
  \item ¿Puede una CNN aprender señal visual útil en un simulador QRS/ST controlado con presupuestos \(K\) muy pequeños?
  \item ¿Se mantiene ese aprendizaje al pasar a Brugada HUCA, un conjunto real con variabilidad clínica y etiqueta binaria de caso confirmado frente a control?
  \item ¿La adaptación de dominio reduce de forma suficiente la brecha entre el dominio sintético y el dominio real?
\end{enumerate}

La parte VLM/ICL queda definida como siguiente fase experimental. No se reportan métricas VLM porque todavía no existen inferencias completas y auditadas. Incluir números no ejecutados rompería la trazabilidad del estudio.

\section{Hipótesis}

El trabajo parte de cinco hipótesis operativas:

\begin{enumerate}
  \item El simulador QRS/ST genera diferencias visuales suficientes para que una CNN aprenda por encima del azar.
  \item El rendimiento CNN mejora al aumentar \(K\) en el dominio sintético, donde las etiquetas proceden de reglas controladas.
  \item La misma CNN pierde rendimiento en Brugada HUCA por cambio de dominio visual y por cambio de significado de la etiqueta.
  \item La adaptación de dominio mejora parcialmente el resultado real, pero no elimina la limitación de pocos datos.
  \item Si la CNN no alcanza rendimiento robusto ni con adaptación, queda justificada la comparación posterior con VLM/ICL como estrategia alternativa para la fase de escasez de etiquetas.
\end{enumerate}

\section{Objetivos}

El objetivo general es construir una canalización reproducible que permita comparar estrategias de ECG como imagen en régimen de pocos datos, con orientación de triaje clínico.

Los objetivos específicos son:

\begin{enumerate}
  \item Generar un conjunto sintético QRS/ST con imágenes de V1, V2 y V3 y etiquetas derivadas de reglas morfológicas.
  \item Construir un conjunto Brugada HUCA procesado, separado por paciente y con exclusiones documentadas.
  \item Entrenar ResNet18 mediante validación leave-one-out para \(K=\{2,4,8,16,32\}\) y tres semillas.
  \item Medir exactitud equilibrada, F1, sensibilidad, especificidad, área ROC, precisión media y matrices de confusión.
  \item Comparar rendimiento sintético y real para cuantificar la brecha de dominio.
  \item Revisar, corregir y ejecutar una rama de adaptación de dominio con SSL, CORAL, MMD y DANN.
  \item Interpretar los resultados CNN como línea base para decidir si entrenar modelos visuales específicos es razonable con pocos datos.
  \item Documentar la pipeline VLM/ICL como comparador central pendiente, con prompts, controles y validación de salida.
\end{enumerate}

\section{Contribuciones}

Las contribuciones cerradas son:

\begin{itemize}
  \item un simulador y un conjunto sintético auditado con 100 pacientes y 300 imágenes;
  \item un preprocesado real de Brugada HUCA con 317 pacientes, 951 imágenes y 46 exclusiones;
  \item resultados CNN completos en sintético y HUCA;
  \item resultados de adaptación de dominio para cuatro métodos;
  \item gráficas, tablas, matrices de confusión y ejemplos Grad CAM;
  \item una lectura negativa útil: la CNN aprende en sintético, pero no alcanza rendimiento robusto en HUCA real con pocos datos ni con adaptación;
  \item una pipeline VLM/ICL preparada como siguiente comparador, sin métricas no auditadas;
  \item una memoria orientada a triaje clínico prudente.
\end{itemize}

\section{Alcance}

El trabajo no valida un producto sanitario. No estima utilidad clínica prospectiva ni sustituye revisión por especialista. La etiqueta sintética deriva de reglas internas y la etiqueta real de HUCA representa caso Brugada confirmado frente a control. Estas dos tareas están relacionadas, pero no son equivalentes.

La línea VLM no es decorativa ni secundaria. Forma parte de la pregunta de investigación completa, porque representa una forma distinta de explotar pocos ejemplos. En esta entrega queda como TODO reproducible por una razón de integridad: no existen todavía predicciones completas, métricas agregadas ni auditoría equivalente a la rama CNN.

\section{Estructura de la memoria}

El capítulo \ref{chap:fundamentos} presenta los fundamentos clínicos y técnicos. El capítulo \ref{chap:relacionados} revisa trabajos relacionados. El capítulo \ref{chap:diseno} fija el diseño experimental. El capítulo \ref{chap:simulador} describe el simulador. El capítulo \ref{chap:modelos} detalla la CNN, la adaptación de dominio y la pipeline VLM pendiente. Los capítulos \ref{chap:resultados} y \ref{chap:datos_reales} presentan los resultados sintéticos y reales. El capítulo \ref{chap:discusion} discute las implicaciones para la comparación CNN frente a VLM/ICL y el capítulo \ref{chap:conclusiones} cierra el trabajo.
